\section{Diagnostic Method}
\label{sec:method}

We now turn the interaction-state hypothesis from Section~\ref{sec:theory} into a rollout-based analysis.
The goal is to separate failures in reaching and lifting the cube from failures that occur afterward, when the policy must retain and transport it.

\subsection{Evaluation setting}

We study the state-based version of \ogbench{} \task{} over five held-out evaluation goals, denoted \textit{task1} - \textit{task5}.
Using state observations removes visual representation learning as a confound and focuses the analysis on state-based value and policy learning.
We evaluate five state-based URL methods: \gciql{}, \crl{}, \fb{}, \rldp{}, and \psm{}.
We report interquartile mean (IQM)~\citep{agarwal-neurips21a} success with stratified-bootstrap confidence intervals for aggregate performance.
We then analyze deterministic closed-loop rollouts using the stage-based decomposition below.

\subsection{Rollout-stage decomposition}

We instantiate the sets from Section~\ref{sec:theory} using the simulator-state thresholds in Table~\ref{tab:failure_criteria}.
Each rollout is successful if all criteria are satisfied; otherwise, it is assigned to the first unmet criterion in the ordered sequence.
This converts aggregate success into a stage-specific failure profile.

\begin{table}[htb]
\centering
\caption{Operational criteria for assigning the first unmet stage in state-based rollouts.}
\label{tab:failure_criteria}
\begin{tabular}{ll}
\toprule
\textbf{Stage}  & \textbf{Criterion} \\
\midrule
Reached cube    & Minimum gripper-to-cube distance $<6$ cm at any point in the episode \\ \addlinespace
Lifted cube     & Cube lifted $>3$ cm above the table with the gripper closed for at least $5$ consecutive steps \\ \addlinespace
Final retention & Final cube height above the table $\geq 3$ cm \\ \addlinespace
Goal reached    & Final cube pose within $5$ cm of the goal pose \\
\bottomrule
\end{tabular}
\end{table}

The criteria are computed directly from simulator states and applied identically to every policy.
They do not depend on a method's learned value, representation, score, or model estimate.
The same decomposition can therefore be applied to goal-conditioned, contrastive, predictive, successor-measure, and model-based methods.

\subsection{Cross-method failure profile}

For each method, we report the percentage of episodes assigned to success, reaching failure, lifting failure, final-retention failure, and goal-reaching failure.
This distribution distinguishes methods that fail to reach or lift the cube from methods that lift it but fail to retain it or reach the goal.
We also report the post-lift share of failures:
\begin{equation}
\text{Post-lift share} = \frac{\text{Final-retention fail}+\text{Goal-reaching fail}} {100-\text{Success}}.
\end{equation}
This quantity measures the fraction of failed episodes in which the cube was lifted before the rollout failed.